\documentclass[11pt]{article}
\usepackage{amsmath}
\usepackage{amsfonts}
\usepackage{amsthm}
\usepackage[utf8]{inputenc}
\usepackage[margin=0.75in]{geometry}

\title{CSC111 Winter 2026 Project 1}
\author{Kevin, Selena}
\date{\today}

\begin{document}
\maketitle

\section*{Running the game}
Run the game by executing \texttt{adventure.py}. No additional libraries or installations are required. The game uses only Python standard library modules.

\section*{Game Map}
Game map showing all locations and their connections:

\begin{verbatim}
       6
       |
       3
      /|\
     / | \
    1  7  4
    |     |
    2     8
    |     |
    5    10
\end{verbatim}

Location IDs:
\begin{itemize}
    \item 1: Bahen Center - First Floor
    \item 2: Bahen Center - Elevator Lobby
    \item 3: St. George Street
    \item 4: Bahen Center - Basement
    \item 5: Bahen Center - Third Floor
    \item 6: Robarts Library - First Floor
    \item 7: Coffee Shop
    \item 8: Computer Lab - Locked (puzzle location)
    \item 9: Robarts Library - Circulation Desk
    \item 10: Computer Lab - Unlocked (submission location)
\end{itemize}

Starting location is: 1

\section*{Game solution}
To win the game, you must collect all required items, solve the puzzle to unlock the computer lab, and use all items at the lab. The player wins only when they are at the submission location (Location 10) and have used all items listed in \verb|submission_required_items|. Here is the complete walkthrough:

\begin{verbatim}
take usb_drive
go east
go north
take project_file
go south
go west
go west
go up
take laptop_charger
go down
go east
go south
go west
solve 2511
enter lab
use usb_drive
use project_file
use laptop_charger
\end{verbatim}

\section*{Lose condition(s)}
The player loses the game if they exceed the maximum number of moves (50 moves). This represents running out of time before the 1pm deadline.

List of commands to lose (by exceeding move limit - repeating go east / go west until the game ends - at 50 moves):

\begin{verbatim}
go east
go west
(Repeat this pattern 50 times to exceed the move limit)
\end{verbatim}

Which parts of your code are involved in this functionality:
\begin{itemize}
    \item File: \texttt{adventure.py}
    \item Class: \texttt{AdventureGame}
    \item Method: \texttt{move\_to\_location()} - increments the \texttt{moves} counter each time the player moves
    \item Method: \texttt{check\_lose\_condition()} - checks if \texttt{moves >= max\_moves}
    \item Attribute: \texttt{max\_moves} - set to 50, loaded from \texttt{game\_data.json}
    \item Attribute: \texttt{moves} - tracks the number of moves made by the player
\end{itemize}

\section*{Inventory}

\begin{enumerate}
\item All location IDs that involve items in the game: 1, 5, 6, 7

\item Item data:
\begin{enumerate}
    \item For Item 1:
    \begin{itemize}
    \item Item name: usb\_drive
    \item Item start location ID: 1
    \item Item target location ID: 10
    \end{itemize}
    \item For Item 2:
    \begin{itemize}
    \item Item name: laptop\_charger
    \item Item start location ID: 5
    \item Item target location ID: 10
    \end{itemize}
    \item For Item 3:
    \begin{itemize}
    \item Item name: project\_file
    \item Item start location ID: 6
    \item Item target location ID: 10
    \end{itemize}
    \item For Item 4:
    \begin{itemize}
    \item Item name: coffee
    \item Item start location ID: 7
    \item Item target location ID: 10

    Coffee is consumable and is removed from inventory after being      used.
    \end{itemize}
\end{enumerate}

    \item Exact command(s) that should be used to pick up an item and check inventory:
    \begin{verbatim}
take usb_drive
go east
go north
take project_file
inventory
go south
go west
go west
go up
take laptop_charger
inventory
    \end{verbatim}
    
    \item Which parts of your code (file, class, function/method) are involved in handling the \texttt{inventory} command:
    \begin{itemize}
        \item File: \texttt{adventure.py}
        \item Function: \texttt{handle\_menu\_command()} - handles the "inventory" menu command
        \item Class: \texttt{AdventureGame}
        \item Attribute: \texttt{inventory} - list of item names the player is carrying
        \item Method: \texttt{get\_item\_by\_name()} - retrieves Item object by name to display description
    \end{itemize}
\end{enumerate}

\section*{Score}
\begin{enumerate}

    \item Players earn score by using items at their target locations. Each item has a point value that is added to the score when used correctly. The first location where players can increase their score is location 10 (Computer Lab - Unlocked), after solving the puzzle and entering the lab.

    Exact list of commands leading up to the first score increase:
    \begin{verbatim}
take usb_drive
go east
go south
take coffee
go north
go west
go south
go west
solve 2511
enter lab
score
use coffee
score
    \end{verbatim}

    \item Copy the list you assigned to \texttt{scores\_demo} in the \texttt{simulation.py} file into this section of the report:
    \begin{verbatim}
take usb_drive
go east
go south
take coffee
go north
go west
go south
go west
solve 2511
enter lab
score
use coffee
score
    \end{verbatim}

    \item Which parts of your code (file, class, function/method) are involved in handling the \texttt{score} functionality:
    \begin{itemize}
        \item File: \texttt{adventure.py}
        \item Function: \texttt{handle\_menu\_command()} - handles the "score" menu command to display current score
        \item Class: \texttt{AdventureGame}
        \item Attribute: \texttt{score} - tracks the player's current score
        \item Method: \texttt{use\_item()} - adds \texttt{target\_points} to score when an item is used at its target location
        \item Class: \texttt{Item} (in \texttt{game\_entities.py})
        \item Attribute: \texttt{target\_points} - the number of points awarded for using the item
    \end{itemize}
\end{enumerate}

\section*{Enhancements}
\begin{enumerate}
    \item Logic Puzzle - Keypad Code
    \begin{itemize}
        \item Brief description: A complex logic puzzle that blocks access to the computer lab. The player must solve a mathematical sequence puzzle to unlock a 4-digit keypad code. The puzzle requires understanding a pattern: starting with 2, each subsequent digit follows the pattern: previous digit × 2, then add 1, alternating. The sequence is: 2 → 5 (2×2+1) → 11 (5×2+1) → 23 (11×2+1), but the code uses the first 4 digits: 2, 5, 1, 1 (2511). The player must use the "solve [code]" command to attempt the solution. If correct, the "enter lab" command becomes available.
        
        \item Complexity level: High
        
        \item Reasons: This is a high-complexity enhancement because:
        \begin{itemize}
            \item It required adding puzzle data structure to the Location class (puzzle dictionary with type, hint, solution, and unlocks\_location)
            \item Implementation of puzzle-solving logic in \texttt{AdventureGame.solve\_puzzle()} method that validates solutions and dynamically adds new commands to location.available\_commands
            \item Integration with the game loop to handle "solve" commands and display puzzle hints
            \item The puzzle requires mathematical reasoning and pattern recognition, making it a substantial gameplay element
            \item Modified game\_data.json to include puzzle information for location 8
            \item Updated simulation.py to properly handle puzzle-solving commands in the simulation
            \item Added puzzle hint display in the main game loop
            \item The enhancement affects multiple files and requires careful state management to ensure the unlocked command appears correctly
        \end{itemize}
        
        \item Name the parts of the code which are involved in this enhancement:
        \begin{itemize}
            \item File: \texttt{game\_entities.py}
            \item Class: \texttt{Location}
            \item Attributes: \texttt{puzzle} (Optional[dict]) - stores puzzle information
            \item File: \texttt{game\_data.json}
            \item Location 8 data: includes "puzzle" field with type, hint, solution, and unlocks\_location
            \item File: \texttt{adventure.py}
            \item Class: \texttt{AdventureGame}
            \item Method: \texttt{solve\_puzzle()} - validates solution and unlocks new location command
            \item Function: \texttt{handle\_solve\_command()} - processes solve commands from user input
            \item Main game loop: displays puzzle hints and handles solve commands
        \end{itemize}
        
        \item Copy the list you assigned to \texttt{enhancements\_demo} in the \texttt{simulation.py} file into this section of the report:
        \begin{verbatim}
go south
go west
solve 1234
solve 2511
enter lab
        \end{verbatim}
    \end{itemize}
\end{enumerate}


\end{document}
